\section{Conditional winning positions}\label{sec:templates}\label{sec:calculus}

The purpose of this section is to justify a finite rule that remains
valid when Breaker may play anywhere on the infinite board.
At a finite position let $M,B\subset\Z^2$ be the finite disjoint
sets currently owned by Maker and Breaker. Their complete contents
are retained throughout the argument.

\begin{definition}\label{def:claim}
Let $A\subseteq T\subset\Z^2$ be finite and let $h\ge1$ be an integer.
The triple $(A,T,h)$ is a \emph{valid claim} if, at every position with
Maker to move and with
\[
 A\subseteq M,\qquad B\cap T=\varnothing,
\]
Maker has already completed a target or has a strategy that completes
one within $h$ further actual Maker moves. We call $A$ the required
set, $T$ the envelope, and $h$ the height.
\end{definition}

This definition imposes no restriction on additional Maker cells,
on Breaker cells outside $T$, or on how the position was reached.
In particular, it applies to arbitrary finite disjoint ownership,
without requiring a prescribed difference between $|M|$ and $|B|$.
Only the starting turn is specified.
A \emph{placement} is an affine map $F(x)=Rx+t$ with $R\in\G$ and
$t\in\Z^2$. A valid claim remains valid under the simultaneous
replacement of $A,T$ by $F(A),F(T)$, with the same height: the
bijection $F$ transports a legal strategy, preserves disjointness
and targets, and does not change the number of moves.

The use of finite sufficient conditions and alternative continuations
has a close antecedent in Sieben's proof trees and situations
\cite[Sections~2--3]{Sieben2008}. His situations consist of a Maker
core and a neighborhood without Breaker stones and are stated at
Breaker's turn; the proof trees use an empty intersection of
continuation territories. The rule below works with arbitrary finite
prior ownership and counts every actual fresh Maker claim. Its
quantitative form applies to all three constructions in this article.

\begin{lemma}[Combination rule]\label{lem:combination}\label{thm:composition}
Let $(A_i,T_i,h_i)$, $1\le i\le r$, be valid claims, where
$1\le r<\infty$, and let $p\in\Z^2$. Define
\begin{equation}\label{eq:combination}
 T=\{p\}\cup\bigcup_{i=1}^rT_i,
 \qquad
 A=\left(\bigcup_{i=1}^r A_i\ \cup\
             \bigcap_{i=1}^rT_i\right)\setminus\{p\},
 \qquad h=1+\max_i h_i.
\end{equation}
Then $(A,T,h)$ is valid. Also, at a position with Breaker to move,
if $A\cup\{p\}\subseteq M$ and $B\cap T=\varnothing$, then Maker
has already won or can win within $h-1$ further Maker moves.
\end{lemma}

\begin{proof}
Every $A_i$ is contained in $T_i$, so $A\subseteq T$.
Suppose first that Maker is to move, $A\subseteq M$,
$B\cap T=\varnothing$, and no target is yet complete.
Since $p\in T$, Breaker does not own $p$.
If $p$ is free, Maker claims it. If $p\in M$, Maker instead claims
any free cell. Such a cell exists because the board is infinite
and $M\cup B$ is finite. This replacement is one real move.
In both cases the resulting Maker set $M'$ contains
\begin{equation}\label{eq:intersection-owned}
 \bigcup_i A_i\ \cup\ \bigcap_iT_i
 \subseteq A\cup\{p\}\subseteq M'.
\end{equation}
If this move completes a target, stop immediately.
Otherwise let $b$ be the next legal Breaker cell. As $b\notin M'$,
\eqref{eq:intersection-owned} implies $b\notin\bigcap_iT_i$.
Choose an index $i$ with $b\notin T_i$.
The older Breaker cells also avoid $T_i$, since $T_i\subseteq T$.
Thus, at the next Maker turn,
\[
 A_i\subseteq M',\qquad (B\cup\{b\})\cap T_i=\varnothing.
\]
The child claim gives a win within $h_i$ more Maker moves.
Including the move already spent gives at most $1+h_i\le h$.

For the Breaker-turn assertion, the inclusions in
\eqref{eq:intersection-owned} already hold with $M$ in place of
$M'$. After Breaker's next move the same surviving child applies,
without spending the parent Maker move. Its bound is at most
$\max_i h_i=h-1$.
\end{proof}

The replacement move cannot invalidate this argument: it only adds
Maker ownership, and its actual ensuing Breaker reply is included
in the proof. It must not be omitted from the count.
A reply outside $T$ misses every child envelope, so distant play
requires no separate case analysis. The Breaker-turn assertion
also shows that $p$ need not have been Maker's most recent claim.

\begin{lemma}[Six bases]\label{lem:bases}
Write the points of $S$ in the order in \eqref{eq:target} as
$s_0,\ldots,s_5$. The claims
\begin{equation}\label{eq:bases}
 (A_j,T_j,h_j)=(S\setminus\{s_j\},S,1),
 \qquad 0\le j<6,
\end{equation}
are valid.
\end{lemma}

\begin{proof}
If Maker owns $s_j$, it already owns all of $S$.
Otherwise $s_j$ is free because Breaker owns none of $S$.
Claiming it completes the target in one move.
\end{proof}

Every application of Lemma~\ref{lem:combination} may use placed
versions of earlier claims. We will describe a finite collection of
such applications and calculate its last required set and height.
The semantic proof is already complete: once that required set is
empty, the final height bounds play from the empty board.
